\documentclass{article}
\usepackage[UTF8]{ctex}
\usepackage{amsmath,amssymb,amsfonts}
\usepackage{enumitem}
\usepackage{geometry}
\geometry{a4paper, margin=1in}
% === 错题本排版优化（不参与题目标识） ===
\setlength{\parindent}{0pt}
\setlength{\parskip}{0.35em plus 0.10em minus 0.08em}
\setlength{\emergencystretch}{3em}
\setlist{itemsep=0.25em, topsep=0.40em, parsep=0pt}
\allowdisplaybreaks[2]
% === 错题本排版优化结束 ===
\begin{document}
\section*{中国石油大学(华东)2024年数学分析真题}
\section*{28．中国石油大学（华东）2024年数学分析试卷}

\subsection{题目1}
1．（10 分）求极限

\[
\lim _{n \rightarrow \infty}\left(\frac{\pi}{\sqrt{n^{2}-1}}+\frac{\pi}{\sqrt{n^{2}-2}}+\cdots+\frac{\pi}{\sqrt{n^{2}-2 n}}\right) .
\]

\subsection{题目2}
2．（10 分）求极限 \(\lim _{x \rightarrow 0}\left(\frac{\sin x}{x}\right)^{\frac{1}{x^{2}}}\) ．

\subsection{题目3}
3．（10 分）对任意的自然数 \(n\) ，有 \(a_{n} \leq b_{n} \leq c_{n}\) ，级数 \(\sum_{n=1}^{\infty} a_{n}\) 与 \(\sum_{n=1}^{\infty} c_{n}\) 均收敛，证明：级数 \(\sum_{n=1}^{\infty} b_{n}\) 也收敛。

\subsection{题目4}
4．（10 分）设 \(f(x)\) 在 \([0,1]\) 上连续，且 \(f(\sin x)+f(\cos x) \neq 0\) ，求积分

\[
\int_{0}^{\frac{\pi}{2}} \frac{f(\sin x)}{f(\sin x)+f(\cos x)} \mathrm{d} x .
\]

\subsection{题目5}
5．（10 分）求 \(x^{2}+(y-b)^{2}=a^{2}(b>a>0)\) 绕 \(x\) 轴旋转所得的旋转体的体积．

\subsection{题目6}
6．（10 分）设数列 \(\left\{a_{n}\right\}\) 满足：存在正数 \(M>0\) ，对一切 \(n\) 有

\[
A_{n}=\left|a_{2}-a_{1}\right|+\left|a_{3}-a_{2}\right|+\cdots+\left|a_{n}-a_{n-1}\right| \leq M .
\]

证明：数列 \(\left\{a_{n}\right\}\) 与 \(\left\{A_{n}\right\}\) 都收敛。

\subsection{题目7}
7．（10 分）记

\[
F(t)=\iiint_{V} x \ln \left(1+x^{2}+y^{2}+z^{2}\right) \mathrm{d} x \mathrm{~d} y \mathrm{~d} z .
\]

其中 \(V: x^{2}+y^{2}+z^{2} \leq t, \sqrt{y^{2}+z^{2}} \leq x\) ，求 \(\frac{\mathrm{d} F(t)}{\mathrm{d} t}\) ．

\subsection{题目8}
8．（10 分）设 \(\Omega\) 是由 \(z=\sqrt{9-x^{2}-y^{2}}\) 与 \(z=\sqrt{x^{2}+y^{2}}\) 围成的区域，\(f(x, y, z)\) 为 \(\Omega\) 上的连续函数，将\\
\(\iiint_{\Omega} f(x, y, z) \mathrm{d} x \mathrm{~d} y \mathrm{~d} z\) 分别用直角坐标系，柱面坐标，球面坐标化为三重积分（无需计算）．

\subsection{题目9}
9．（10 分）设 \(f(x)\) 在 \(x=0\) 的某邻域内存在连续的二阶导数，且 \(\lim _{x \rightarrow 0} \frac{f(x)}{x}=0\) ．证明：级数 \(\sum_{n=1}^{\infty} f\left(\frac{1}{n}\right)\)绝对收敛。

\subsection{题目10}
10．（10 分）计算曲面积分

\[
\iint_{S}(x-x y) \mathrm{d} y \mathrm{~d} z+(y-y z) \mathrm{d} z \mathrm{~d} x+(z-x y) \mathrm{d} x \mathrm{~d} y .
\]

其中 \(S\) 是以原点为球心，半径为 \(a\) 的上半球面，方向取上侧。

\subsection{题目11}
11．（10 分）求 \(y=\ln \left(1-3 x+x^{2}\right)\) 的麦克劳林级数展开式．

\subsection{题目12}
12．（10 分）设 \(f(x)\) 在 \((-\infty,+\infty)\) 上连续，且 \(\lim _{x \rightarrow \infty} \frac{f(x)}{x}=0\) ，证明：存在 \(\xi \in(-\infty,+\infty)\) ，使得\(f(\xi)=\xi\) ．

\subsection{题目13}
13．（10 分）设 \(f(x, y)=x+y+x y\) ，曲线 \(C: x^{2}+y^{2}=2\) 。

\begin{itemize}
  \item[（1）] 求 \(f(x)\) 在 \(P(1,1)\) 沿方向 \(1=(-1,-1)\) 的方向导数．
  \item[（2）] 求 \(f(x)\) 在曲线 \(C\) 上的最大方向导数．
\end{itemize}

\subsection{题目14}
14．（10 分）设 \(f(x)\) 在 \([1,2]\) 上可导，且 \(f^{\prime}(x) \neq 1, \forall x \in[1,2]\) ．若对任意的 \(x \in[1,2]\) ，有 \(f(x) \in(1,2)\) ，证明：存在唯一的 \(\xi \in[1,2]\) ，使得 \(f(\xi)=\xi\) 。

\subsection{题目15}
15．（10 分）设 \(f(x)\) 在 \([a, b]\) 上可积，证明：\(f^{2}(x)\) 在 \([a, b]\) 上也可积．
\end{document}
